\documentclass[12pt]{article}
\usepackage[utf8]{inputenc}
\usepackage[T1]{fontenc}
\usepackage{hyperref}
\title{Southwestern Fars Language Resources for K-12 Curriculum Development}
\author{Compiled by Claude Sonnet 5.5 (claude-sonnet-5-5) for Andr\'as Kornai}
\date{}
\begin{document}
\maketitle

\section{Introduction}
Southwestern Fars (ISO 639-3 \texttt{fay}, Glottocode sout2645) is a group of Southwestern Iranian varieties in Fars province, Iran, with roughly 100,000 speakers according to Global Recordings Network. Glottolog links it to the Lari group; Global Recordings lists Basseri, Delvari and Old Shirazi as varieties. The searches found only catalogue and encyclopedic entries. No textbooks, primers or school programmes were found, and this dossier is accordingly short.

\section{Primary Educational Resources}
None found.

\section{Linguistic and Reference Materials}
\begin{itemize}
\item Glottolog entry for Southwestern Fars: \url{https://glottolog.org/resource/languoid/id/sout2645}
\item ``Dialects of Fars'' (Wikipedia): \url{https://en.wikipedia.org/wiki/Dialects_of_Fars}
\item ``Old Kazeruni dialect'' (Wikipedia), an extinct neighbouring Southwestern Iranian variety, which per the search summary has no ISO code of its own: \url{https://en.wikipedia.org/wiki/Old_Kazeruni_dialect}
\end{itemize}

\section{Religious and Community Resources}
None found beyond the Global Recordings listing below.

\section{Digital and Audio Resources}
\begin{itemize}
\item Global Recordings Network page for Southwestern Fars (gospel audio recordings catalogue): \url{https://globalrecordings.net/en/language/fay}
\end{itemize}

\section{Curriculum Development Context}
Persian is the language of instruction in Iran, and no school use of Southwestern Fars was found. Any curriculum would need to start from fieldwork documentation; the sources above give only classification and variety names.

\end{document}